\documentclass[10pt,a4paper]{amsart}
\usepackage[margin=19mm]{geometry}
\usepackage{amsmath,amssymb,mathtools}
\usepackage[scr=rsfs,cal=euler]{mathalfa}
\usepackage{xfrac}
\usepackage[unicode,hidelinks,bookmarks=false]{hyperref}
\newcommand{\ie}{i.e.\ }
\newcommand{\cf}{cf.\ }
\newcommand{\resp}{respectively\ }
\newcommand{\Subsection}{\subsection}
\providecommand{\nobreakdash}{\nobreak\hbox{-}}
\setlength{\emergencystretch}{2em}
\multlinegap0pt
\usepackage[shortlabels]{enumitem}
\usepackage{amssymb}
\usepackage{mathtools}
\usepackage{xcolor}
\newtheorem{lem}{Lemma}
\newtheorem{theorem}{Theorem}
\newcommand{\bC}{\mathbb{C}}
\newcommand{\mfX}{\mathfrak{X}}
\newcommand{\mfg}{\mathfrak{g}}
\title{Stratified Pseudobundles and Quantization}
\author{Ethan Ross}
\date{Source: arXiv:2601.03544v1 (CC BY 4.0)}
\begin{document}
\maketitle

\noindent\emph{Source and licence.} Stratified Pseudobundles and Quantization. Version arXiv:2601.03544v1 (CC BY 4.0), distributed by arXiv under the Creative Commons Attribution 4.0 International licence, http://creativecommons.org/licenses/by/4.0/, as recorded in the arXiv licence metadata for this exact version and shown on the abstract page https://arxiv.org/abs/2601.03544v1. Full licence notice: \texttt{licenses/arxiv-2601.03544-CC-BY-4.0.txt}.\par\smallskip

\noindent\textit{Editorial scope.} Counted unit: the article's theorem non-singular (source0.tex lines 4451--4479). The article's complete original proof of it is delivered with it (source0.tex lines 4481--4560, one \texttt{proof} environment of 80 lines). No mathematics is written, repaired or re-derived: the statement and the proof are the article's own lines, in the article's own notation, and no retained line is substituted or re-flowed.  Retained uncounted as the source-local context the proof needs: lines 4362--4380.  Nothing was omitted from any retained span, so the package carries no omission record.  No retained line contains a citation, so the wrapper carries no bibliography and no bibliography entry.  No retained line contains a figure environment, an image-inclusion command or drawing code, so no image is delivered and the package declares no asset dependency.  Editorial formatting only, in addition: an AMS \texttt{amsart} preamble that restates the article's own declarations and macros used by the retained lines, namely the environments \texttt{lem}, \texttt{theorem} on separate counters, together with the standard packages those lines need.  The retained environments print in this document as Lemma 1, Theorem 1 in order of appearance, whereas the article prints the same items with the numbering of its own full document.  The complete native source of the article as distributed in the arXiv e-print for this exact version is bundled unchanged as \texttt{sources/192232/source0.tex}.
\begin{lem}\label{lem: non-singular reduction of line bundle}
Suppose $J^{-1}(0)\subseteq M$ is an embedded submanifold and, furthermore, suppose $J^{-1}(0)$ has only one orbit-type. Write $\pi_M:J^{-1}(0)\to M_0$ and $\pi_\mathcal{L}:\mathcal{L}|_{J^{-1}(0)}\to \mathcal{L}_0$ for the quotient maps. If the map
\begin{equation}
J^{-1}(0)\times \Gamma(\mathcal{L})^G\to \mathcal{L};\quad (x,\sigma)\mapsto \sigma(x)
\end{equation}
is surjective, then the following holds.
\begin{itemize}
    \item[(i)] The induced map $\mathcal{L}_0\to M_0$ is a smooth complex line bundle over $M_0$.
    \item[(ii)] The projection $\pi_\mathcal{L}:\mathcal{L}|_{J^{-1}(0)}\to \mathcal{L}_0$ induces a bijection
    $$
    \kappa:\Gamma(\mathcal{L}|_{J^{-1}(0)})^G\to \Gamma(\mathcal{L}_0)
    $$
    \item[(iii)] The map
    $$
    \kappa:\Gamma(\mathcal{L})^G\to \Gamma(\mathcal{L}_0);\quad \sigma\mapsto \kappa(\sigma|_{J^{-1}(0)})
    $$
    is a surjective map. 
\end{itemize}
\end{lem}

\begin{theorem}\label{thm: can reduce non-singular}
Suppose $J^{-1}(0)\subseteq M$ is an embedded submanifold with only one orbit type and suppose
$$
J^{-1}(0)\times \Gamma(\mathcal{L})^G\to \mathcal{L};\quad (x,\sigma)\mapsto \sigma(x)
$$
is surjective. Then the following holds.
\begin{itemize}
    \item[(i)] Suppose $\sigma_1,\sigma_2\in\Gamma(\mathcal{L})^G$ and $V_1,V_2\in\mfX(M,J^{-1}(0))^G$ satisfy 
    $$
    \sigma_1|_{J^{-1}(0)}=\sigma_2|_{J^{-1}(0)}
    $$
    and 
    $$
    (\pi_0)_*(V_1|_{J^{-1}(0)})=(\pi_0)_*(V_2|_{J^{-1}(0)}).
    $$
    Then,
    \begin{equation}
    \nabla_{V_1}\sigma_1|_{J^{-1}(0)}=\nabla_{V_2}\sigma_2|_{J^{-1}(0)}.
    \end{equation}
    \item[(ii)] $\mathcal{L}_0=(\mathcal{L}|_{J^{-1}(0)})/G$ is canonically a prequantum line bundle over $M_0=J^{-1}(0)/G$ with prequantum connection $\nabla^0$ given by
    $$
    \nabla^0_V\sigma=\kappa(\nabla_W\tau),
    $$
    where $W\in \mfX(M,J^{-1}(0))^G$ and $\tau\in\Gamma(\mathcal{L})^G$ satisfy
    $$
    (\pi_0)_*W|_{J^{-1}(0)}=V\quad\text{and}\quad \kappa(\tau)=\sigma.
    $$
\end{itemize}
\end{theorem}

\begin{proof}
\begin{itemize}
    \item[(i)] Fix equivariant sections $\sigma_1,\sigma_2\in\Gamma(\mathcal{L})^G$ and restrictable vector fields $V_1,V_2\in\mfX(M,J^{-1}(0))^G$ satisfying $\sigma_1|_{J^{-1}(0)}=\sigma_2|_{J^{-1}(0)}$ and $(\pi_0)_*(V_1|_{J^{-1}(0)})=(\pi_0)_*(V_2|_{J^{-1}(0)})$. Observe that
    $$
    \nabla_{V_1}\sigma_1-\nabla_{V_2}\sigma_2=\nabla_{V_1}(\sigma_1-\sigma_2)+\nabla_{V_1-V_2}\sigma_2
    $$
    I claim both terms vanish after restricting to $J^{-1}(0)$. 

    \

    To show the first term vanishes, fix $x\in J^{-1}(0)$ and let $U\subseteq M$ be a trivialization neighbourhood for $\mathcal{L}$. Then, after choosing an isomorphism $\mathcal{L}|_U\to U\times \bC$, we can identify $\Gamma(\mathcal{L}|_U)$ with $C^\infty(U,\bC)$. Thus, since $\sigma_1=\sigma_2$ over $J^{-1}(0)$, we can identify $\sigma_1-\sigma_2$ with a function $f\in C^\infty(U,\bC)$ with $f|_{J^{-1}(0)}=0$. Furthermore, we can find a local primitive $\alpha\in\Omega^1(U)$ of $\omega|_U$ so that
    $$
    \nabla=d+i \alpha.
    $$
    Hence, 
    $$
    \nabla_{V_1}(f)=V_1f+i \alpha(V_1)f.
    $$
    Using the fact that $f(x)=0$, we get
    $$
    \nabla_{V_1}f(x)=(V_1)_xf.
    $$
    Since $V_1$ is tangent to $J^{-1}(0)$ and $f$ vanishes on $J^{-1}(0)$, we have $(V_1)_xf=0$. Since this holds in any trivializing neighbourhood intersecting $J^{-1}(0)$, we obtain $\nabla_{V_1}(\sigma_1-\sigma_2)|_{J^{-1}(0)}=0$.

    \

    Now for the second term. Since $(\pi_0)_*V_1|_{J^{-1}(0)}=(\pi_0)_*V_2|_{J^{-1}(0)}$, we have
    $$
    (\pi_0)_*((V_1-V_2)|_{J^{-1}(0)})=0.
    $$
    Since $\pi_0:J^{-1}(0)\to M_0$ is the quotient map by the $G$-action and the action of $G$ on $J^{-1}(0)$ is regular, there exists functions $h_1,\dots,h_M\in C^\infty(J^{-1}(0))$ and Lie algebra elements $\xi_1,\dots,\xi_M\in\mfg$ so that
    $$
    (V_1-V_2)|_{J^{-1}(0)}=\sum_{i=1}^M h_i (\xi_i)_{J^{-1}(0)},
    $$
    where $(\xi_i)_{J^{-1}(0)}\in \mfX(J^{-1}(0))$ is the fundamental vector field associated to $\xi_i$. Thus,
    $$
    \nabla_{V_1-V_2}\sigma_2|_{J^{-1}(0)}=\sum_{i=1}^Mh_i\nabla_{(\xi_i)_M}\sigma_2|_{J^{-1}(0)}.
    $$
    By the equivariance of $\nabla$ and $\sigma_2$, we have
    $$
    \nabla_{(\xi_i)_M}\sigma_2=\xi_i\cdot \sigma+i J^{\xi_i}\sigma_2=+i J^{\xi_i}\sigma_2
    $$
    Thus, restricting to $J^{-1}(0)$ we have
    $$
    \nabla_{V_1-V_2}\sigma_2|_{J^{-1}(0)}=0.
    $$

    \item[(ii)] $\mathcal{L}_0$ is a complex line bundle is statement (i) from Lemma \ref{lem: non-singular reduction of line bundle}. Now let us show $\nabla^0$ is well-defined. Fix $V\in\mfX(M_0)$ and $\sigma\in\Gamma(\mathcal{L}_0)$ and suppose $W_1,W_2\in\mfX(M,J^{-1}(0))^G$ are two lifts of $V$ and $\tau_1,\tau_2\in \Gamma(\mathcal{L})^G$ are two lifts of $\sigma$. Then, by construction we have
    $$
    (\pi_0)_*W_1|_{J^{-1}(0)}=(\pi_0)_*W_2|_{J^{-1}(0)}=V
    $$
    and
    $$
    \kappa(\tau_1)=\kappa(\tau_2)=\sigma.
    $$
    Hence, by (i)
    $$
    \nabla_{W_1}\tau_1|_{J^{-1}(0)}=\nabla_{W_2}\tau_2|_{J^{-1}(0)}.
    $$
    Since both $\nabla_{W_1}\tau_1$ and $\nabla_{W_2}\tau_2$ are equivariant sections of $\mathcal{L}$, we conclude by (iii) of Lemma \ref{lem: non-singular reduction of line bundle} that
    $$
    \kappa(\nabla_{W_1}\tau_1)=\kappa(\nabla_{W_2}\tau_2).
    $$
    Hence, $\nabla^0$ is well-defined. $\nabla^0$ is also clearly smooth since the image of $\kappa$ is the smooth sections of $\Gamma(\mathcal{L}_0)$. Finally, we show $\nabla^0$ is prequantum. Fix $V_1,V_2\in\mfX(M_0)$ with lifts $W_1,W_2\in\mfX(M,J^{-1}(0))^G$ and a section $\sigma\in\Gamma(\mathcal{L}_0)$ with a lift $\tau\in\Gamma(\mathcal{L})^G$. Unfolding definitions shows
    $$
    [\nabla^0_{V_1},\nabla^0_{V_2}]\sigma=\kappa([\nabla_{W_1},\nabla_{W_2}]\tau)
    $$
    and since the map $(\pi_0)_*:\mfX(M,J^{-1}(0))^G\to \mfX(M_0)$ is a Lie algebra homomorphism,
    $$
    \nabla^0_{[V_1,V_2]}\sigma=\kappa(\nabla_{[W_1,W_2]}\tau).
    $$
    Hence,
    \begin{align*}
    [\nabla^0_{V_1},\nabla^0_{V_2}]\sigma-\nabla^0_{[V_1,V_2]}\sigma&=\kappa([\nabla_{W_1},\nabla_{W_2}]\tau-\nabla_{[W_1,W_2]}\tau)\\
    &=\kappa(-i \omega(W_1,W_2)\tau)\\
    &=-i \omega_0(V_1,V_2)\sigma.
    \end{align*}
    where we used the fact that $\nabla$ is prequantum and that $\omega|_{J^{-1}(0)}=\pi_0^*\omega_0$. Thus, $\nabla^0$ is prequantum.
\end{itemize}
\end{proof}

\end{document}
